\section{Phần A --- Cài đặt từ đầu}
\label{sec:phan-a}

Đề yêu cầu Phần~A phải \emph{dẫn xuất từ bài toán tối ưu hoặc mô hình xác suất},
cài đặt từ đầu bằng NumPy hoặc PyTorch mà không dùng hàm huấn luyện dựng sẵn, rồi
đối chiếu với thư viện chuẩn \cite[tr.~1, mục I]{de}. Mục này theo đúng ba bước đó
cho từng hạng mục. Thư viện chuẩn chỉ xuất hiện ở phía đối chiếu, không xuất hiện
trong đường tính toán của nhóm.

Mã dùng lại nằm trong \texttt{src/}; notebook điều khiển và giải thích nằm trong
\texttt{notebooks/}. Toàn bộ phép kiểm ở mục này chạy trên CPU cục bộ.

%--------------------------------------------------------------------
\subsection{A.1 --- Các thành phần của khối decoder}
\label{sec:a1}

\subsubsection{Dẫn xuất: attention như một phép lấy trung bình có trọng số}

Xuất phát từ bài toán truy xuất mềm. Cho một truy vấn $q\in\R^{d_k}$ và $n$ cặp
khóa--giá trị $(k_j,v_j)$, ta muốn lấy ra giá trị ứng với khóa ``giống'' truy vấn
nhất. Phép lấy cứng $v_{\argmax_j\langle q,k_j\rangle}$ không khả vi, nên thay bằng
một phân phối trên các khóa rồi lấy kỳ vọng. Chọn họ phân phối Gibbs với thế năng
là tích vô hướng:
\[
  \alpha_j=\frac{\exp(\langle q,k_j\rangle/T)}{\sum_{a=1}^{n}\exp(\langle q,k_a\rangle/T)},
  \qquad \text{đầu ra} = \sum_{j=1}^{n}\alpha_j v_j .
\]
Đây chính là nghiệm của bài toán tối ưu có phạt entropy
\[
  \max_{\alpha\in\Delta^{n-1}}\ \sum_j \alpha_j\langle q,k_j\rangle
    + T\Bigl(-\sum_j\alpha_j\log\alpha_j\Bigr),
\]
với $\Delta^{n-1}$ là simplex xác suất; nghiệm là softmax của $\langle q,k_j\rangle/T$.
Cấu trúc này lặp lại ở Mục~\ref{sec:b5} với $T=\beta$ và thế năng là reward, nên
attention và DPO dùng cùng một dạng nghiệm Gibbs.

Chọn $T=\sqrt{d_k}$ không phải tùy ý. Dưới giả thiết các thành phần của $q,k$ độc
lập, kỳ vọng $0$, phương sai $1$, ta chứng minh ở Mục~\ref{sec:b1} rằng
$\Var(q^\top k)=d_k$; chia cho $\sqrt{d_k}$ đưa phương sai mỗi điểm về $1$ và thu
nhỏ mọi chênh lệch điểm theo cùng hệ số, nhờ đó giảm bão hòa softmax
\cite[§3.2.1]{attention}.

Viết theo ma trận cho cả $n$ truy vấn:
\begin{equation}
  \Attn(Q,K,V)=\softmax\!\left(\frac{QK^\top}{\sqrt{d_k}}\right)V,
  \label{eq:sdpa}
\end{equation}
đúng phương trình (1) của \cite{attention}.

\subsubsection{Mặt nạ nhân quả như một ràng buộc hỗ trợ}

Mô hình ngôn ngữ nhân quả phân rã xác suất chuỗi theo quy tắc dây chuyền:
\begin{equation}
  p_\theta(y_1,\ldots,y_T)=\prod_{t=1}^{T}p_\theta(y_t\mid y_{<t}).
  \label{eq:chain}
\end{equation}
Để \eqref{eq:chain} đúng, biểu diễn tại vị trí $t$ không được phụ thuộc
$y_{>t}$. Ràng buộc này áp vào \eqref{eq:sdpa} bằng cách buộc $\alpha_{tj}=0$ với
$j>t$, tương đương cộng $-\infty$ vào điểm số trước softmax. Vi phạm ràng buộc này
làm \eqref{eq:chain} sai, nên hàm mất mát không còn là log-hợp lý của mô hình nào
cả --- đây là lý do mặt nạ là vấn đề đúng/sai, không phải lựa chọn.

Cài đặt đặt mặt nạ dưới dạng boolean ``bị chặn'', giữ nguyên đường chéo:

\begin{quote}\small\ttfamily
def causal\_mask(length, *, device=None):\\
\hspace*{2em}return torch.ones(length, length, dtype=torch.bool, device=device).triu\_(1)
\end{quote}

Hàm \texttt{scaled\_dot\_product\_attention} nhận mặt nạ này và dùng
\texttt{masked\_fill\_} với $-\infty$ trước \texttt{softmax}. Mỗi truy vấn luôn còn
ít nhất một khóa hợp lệ (chính nó), nên mẫu số softmax dương và không sinh
\texttt{NaN}. Tham số \texttt{blocked=None} nghĩa là \emph{không che}, không phải
chế độ nhân quả --- phân biệt này được ghi trong docstring để tránh dùng sai.

\subsubsection{Chú ý nhiều đầu}

Một đầu attention chỉ tạo được một phân phối trọng số trên mỗi vị trí. Để mô hình
dùng nhiều quan hệ song song, chia $\dmodel$ thành $h$ đầu, mỗi đầu có
$d_h=\dmodel/h$, tính attention riêng rồi ghép và chiếu
\cite[§3.2.2]{attention}:
\[
  \mathrm{MHA}(X)=\mathrm{Concat}(H_1,\ldots,H_h)W_O,\qquad
  H_i=\Attn(XW_Q^{(i)},XW_K^{(i)},XW_V^{(i)}).
\]
Vì $h\cdot d_h=\dmodel$, tổng chi phí các phép chiếu không đổi so với một đầu có
cùng $\dmodel$ --- một điểm thường bị hỏi ở vấn đáp và được dùng lại ở
Mục~\ref{sec:c2}.

Cài đặt chiếu Q/K/V \emph{một lần} bằng \texttt{nn.Linear(d\_model, 3*d\_model)} rồi
tách, thay vì ba lớp riêng. Phép biến đổi hình dạng giữ đúng thứ tự trục:

\begin{quote}\small\ttfamily
q, k, v = self.qkv(x).view(batch, length, 3, self.n\_heads, self.head\_dim)\\
\hspace*{6em}.permute(2, 0, 3, 1, 4).unbind(0)
\end{quote}

Cho \texttt{[B,T,3,H,Dh]} $\to$ ba tensor \texttt{[B,H,T,Dh]}. Trục đầu và trục
token không bị trộn; nếu trộn, mỗi đầu sẽ nhận một lát đặc trưng sai và kết quả vẫn
có hình dạng đúng nhưng sai về ngữ nghĩa. Đây là lỗi im lặng nên notebook A.1 kiểm
riêng bằng cách so với \texttt{nn.MultiheadAttention} dùng \emph{cùng trọng số}.

\subsubsection{Mã hóa vị trí hình sin}

Self-attention không có mã hóa vị trí là đẳng biến hoán vị (chứng minh ở
Mục~\ref{sec:b2}), nên không phân biệt được thứ tự token. Transformer gốc cộng một
bảng cố định \cite[§3.5]{attention}:
\[
  \PE(p)_{2i}=\sin\!\left(p\,\omega_i\right),\quad
  \PE(p)_{2i+1}=\cos\!\left(p\,\omega_i\right),\quad
  \omega_i=10000^{-2i/\dmodel}.
\]
Tần số được tính trong miền log để tránh tràn số:

\begin{quote}\small\ttfamily
freq = torch.exp(torch.arange(0, d\_model, 2, dtype=torch.float32)\\
\hspace*{6em}* (-math.log(10000.0) / d\_model))
\end{quote}

Bảng được đăng ký bằng \texttt{register\_buffer(\ldots, persistent=False)}: nó là
hằng số, không phải tham số học, và không cần lưu vào \texttt{state\_dict} vì tái
tạo được từ cấu hình. PE được cộng \emph{đúng một lần} trước chồng block.

Cài đặt chấp nhận $\dmodel$ lẻ; khi đó tọa độ $\sin$ cuối không có $\cos$ đi kèm.
Mục~\ref{sec:b3} chứng minh rằng \emph{tính chất dịch vị trí không còn đúng} trong
trường hợp đó. Thí nghiệm A.2 dùng $\dmodel=64$ nên không chạm phải giới hạn này,
nhưng giới hạn được ghi rõ thay vì bỏ qua.

\subsubsection{Khối decoder pre-LN}

Khối ghép hai nhánh residual, chuẩn hóa \emph{trước} mỗi nhánh:
\[
  x \leftarrow x + \mathrm{Attn}(\mathrm{LN}_1(x)),\qquad
  x \leftarrow x + \mathrm{FFN}(\mathrm{LN}_2(x)).
\]
Nhóm chọn pre-LN thay vì post-LN của bài gốc vì gradient đi qua nhánh residual
không bị LayerNorm chắn, nên huấn luyện ổn định hơn ở learning rate cố định mà
không cần warmup. Đây là lựa chọn thiết kế của nhóm, được ghi nhận là khác bài gốc.

FFN theo từng vị trí, hai lớp tuyến tính với ReLU ở giữa
\cite[§3.3]{attention}. Không có cross-attention vì mô hình chỉ có decoder.

\subsubsection{Đối chiếu \dalam}

\begin{center}
\begin{tabular}{p{0.26\textwidth}p{0.30\textwidth}p{0.34\textwidth}}
\toprule
\textbf{Hạng mục} & \textbf{Đối chiếu với} & \textbf{Kết quả}\\
\midrule
A.1a attention & \texttt{F.scaled\_dot\_product\_}\allowbreak\texttt{attention} &
Khớp cả khi có và không có mặt nạ\\
A.1b nhiều đầu & \texttt{nn.MultiheadAttention} cùng trọng số &
Khớp đầu ra \emph{và} gradient theo đầu vào\\
A.1c mặt nạ & Trung bình tiền tố tính tay &
Đạt; kiểm thêm đổi hậu tố và cắt ngắn chuỗi\\
A.1d PE & \texttt{sin}/\texttt{cos} vô hướng &
Khớp bảng đã cộng vào embedding, kể cả $\dmodel$ lẻ\\
A.1e decoder & Chạy tiền tố trên CPU &
Luồng mặt nạ và các thành phần tự cài được dùng trực tiếp\\
\bottomrule
\end{tabular}
\end{center}

Phép kiểm \emph{đổi hậu tố} đáng chú ý: giữ nguyên $y_{\le t}$ và thay đổi
$y_{>t}$, đầu ra tại vị trí $t$ phải không đổi. Phép kiểm này bắt được rò rỉ tương
lai mà phép so sánh hình dạng không bắt được.

%--------------------------------------------------------------------
\subsection{A.2 --- Mô hình ngôn ngữ ký tự và perplexity}
\label{sec:a2}

\subsubsection{Dẫn xuất: hợp lý cực đại trên chuỗi ký tự}

Mục tiêu là cực đại log-hợp lý của dữ liệu dưới \eqref{eq:chain}, tương đương cực
tiểu cross-entropy trung bình theo token. Mục~\ref{sec:b6} chứng minh mục tiêu này
chỉ khác $\KL(\pdata\|p_\theta)$ một hằng số không phụ thuộc $\theta$, nên hai cách
phát biểu cho cùng tập nghiệm. Perplexity là hàm mũ của cross-entropy trung bình
theo token, tức nghịch đảo trung bình nhân của xác suất token đúng.

\subsubsection{Dữ liệu và quyền sử dụng}

Corpus là bản \emph{Truyện Kiều} do Trương Vĩnh Ký phiên âm, in năm 1911, thuộc
phạm vi công cộng theo thời hạn bảo hộ. Nhóm chuyển sang bản này sau khi rà soát
cho thấy căn cứ quyền đối với bản số hóa trước đó (revision \texttt{206607}) chưa
đủ. Toàn bộ quá trình rà soát, kể cả kết luận ``chưa đủ căn cứ'', được giữ trong
\texttt{docs/data/char-corpus.md} thay vì xóa đi.

\begin{center}
\begin{tabular}{lr}
\toprule
\textbf{Chỉ số} & \textbf{Giá trị}\\
\midrule
Số dòng thơ & 3.254\\
Số code point NFC cả corpus & 104.181\\
Tách train (prefix 90\%) & 93.762 code point\\
Tách validation (suffix 10\%) & 10.419 code point\\
Kích thước từ vựng & 127 (gồm UNK, chỉ học từ train)\\
Độ dài ngữ cảnh / stride & 128 / 128\\
Số cửa sổ train / validation & 732 / 81\\
Code point đuôi không dùng (train / validation) & 65 / 50\\
Target được chấm ở validation / số UNK & 10.368 / 1\\
\bottomrule
\end{tabular}
\end{center}

Ba quyết định về dữ liệu đáng nêu, vì mỗi quyết định chặn một đường rò rỉ:

\begin{enumerate}[leftmargin=2em]
\item \textbf{Tách train/validation \emph{trước} khi học từ vựng.} Nếu học từ vựng
trên toàn corpus, thông tin về ký tự chỉ xuất hiện ở validation sẽ lọt vào mô hình
qua bảng embedding. Từ vựng được học chỉ từ train; ký tự validation chưa gặp nhận
ID \texttt{<UNK>} bằng $0$.
\item \textbf{Tách trước khi tạo cửa sổ.} Nếu tạo cửa sổ trước rồi chia, một cửa sổ
có thể bắc qua ranh giới và đưa ký tự validation vào đầu vào train.
\item \textbf{\texttt{<UNK>} vẫn được chấm.} Có đúng một đích \texttt{<UNK>} trong
validation; nó không bị loại khỏi mẫu số perplexity. Loại nó ra sẽ làm PPL tốt hơn
một cách giả tạo.
\end{enumerate}

\subsubsection{Tổng hợp perplexity}

Quy tắc tổng hợp quan trọng ngang với công thức: cộng tổng log-xác suất âm rồi chia
tổng số token được chấm, \emph{không} lấy trung bình các PPL từng batch
(Mục~\ref{sec:b6} giải thích vì sao). Cài đặt tích lũy ở \texttt{float64} với
\texttt{reduction='sum'}:

\begin{quote}\small\ttfamily
nll\_sum += F.cross\_entropy(logits.flatten(0, 1), batch[:, 1:].reshape(-1),\\
\hspace*{6em}reduction='sum')\\
token\_count = windows.size(0) * (windows.size(1) - 1)\\
nll = nll\_sum.item() / token\_count\\
\ldots 'validation\_ppl': math.exp(nll)
\end{quote}

\subsubsection{Kết quả \dalam}

Năm seed độc lập, mỗi seed 5 epoch với 2 decoder block, $\dmodel=64$, 4 đầu,
$d_{ff}=128$, batch size 16, AdamW learning rate $0{,}001$. Các seed dùng cùng một
split; mỗi lượt chấm 10.368 ký tự đích (81 cửa sổ $\times$ 128).

\begin{center}
\begin{tabular}{rrr}
\toprule
\textbf{Seed} & \textbf{NLL validation (nat/ký tự)} & \textbf{PPL validation}\\
\midrule
41 & 2,326985 & 10,2470\\
42 & 2,328789 & 10,2655\\
43 & 2,321075 & 10,1866\\
44 & 2,318854 & 10,1640\\
45 & 2,313612 & 10,1109\\
\midrule
\multicolumn{2}{r}{\textbf{Trung bình $\pm$ SD mẫu}} & $\mathbf{10{,}1948 \pm 0{,}0628}$\\
\bottomrule
\end{tabular}
\end{center}

Con số $\pm 0{,}0628$ là độ lệch chuẩn mẫu (mẫu số $4$) của 5 lần huấn luyện độc
lập \emph{trên cùng một split}. Nó đo độ biến thiên do khởi tạo và thứ tự trộn,
không phải khoảng tin cậy, và không phải đánh giá trên tập test độc lập. Việc này
thỏa yêu cầu tái lập ở phần~II của đề (ít nhất 5 lần chạy, báo cáo trung bình $\pm$
độ lệch chuẩn) cho riêng hạng mục A.2.

\textbf{Giới hạn cần nêu.} Huấn luyện chạy trên Kaggle, phiên hiển thị cấu hình
\texttt{GPU T4 x2}; đây là cấu hình nền tảng hiển thị, không phải xác nhận mô hình
đã dùng cả hai GPU. Tệp \texttt{summary.json} nhận được không chứa phiên bản
PyTorch/CUDA và không kèm checkpoint, nên nhóm \emph{chưa} đối chiếu được log môi
trường. Mục~\ref{sec:tai-lap} ghi việc này là hạng mục còn thiếu.

%--------------------------------------------------------------------
\subsection{A.3 --- Lớp LoRA có hợp nhất trọng số}
\label{sec:a3}

\subsubsection{Dẫn xuất: giới hạn cập nhật vào một không gian hạng thấp}

Tinh chỉnh toàn phần cập nhật cả $W\in\R^{d\times k}$, tức $dk$ tham số cho mỗi
lớp. LoRA ràng buộc cập nhật nằm trong tập các ma trận hạng không quá $r$
\cite[§4.1]{lora}:
\[
  W'=W+\Delta W,\qquad \Delta W=\frac{\alpha}{r}BA,\qquad
  B\in\R^{d\times r},\ A\in\R^{r\times k}.
\]
Đây là bài toán tối ưu trên một đa tạp tham số nhỏ hơn hẳn: thay vì $dk$ chiều,
chỉ còn $r(d+k)$ chiều (Mục~\ref{sec:b4}). Hệ số $\alpha/r$ giữ độ lớn cập nhật
tương đối ổn định khi đổi $r$, nên learning rate không phải điều chỉnh lại theo
từng rank.

\subsubsection{Cài đặt}

Lớp bọc trực tiếp \texttt{nn.Linear} nền và đóng băng \emph{cả weight lẫn bias}:

\begin{quote}\small\ttfamily
self.scale = alpha / rank\\
self.A = nn.Parameter(base.weight.new\_empty(rank, base.in\_features))\\
self.B = nn.Parameter(base.weight.new\_zeros(base.out\_features, rank))\\
nn.init.normal\_(self.A, std=base.in\_features ** -0.5)\\
self.base = base.requires\_grad\_(False)
\end{quote}

Bốn chi tiết có chủ đích:

\begin{enumerate}[leftmargin=2em]
\item $B$ khởi tạo bằng $0$ nên $W'=W$ tại bước đầu, bất kể $A$
(Mục~\ref{sec:b4}). Mô hình xuất phát đúng bằng mô hình nền.
\item $A$ khởi tạo ngẫu nhiên với $\mathrm{std}=k^{-1/2}$. Nếu khởi tạo \emph{cả
hai} bằng $0$ thì gradient theo cả $A$ và $B$ đều triệt tiêu và adapter không bao
giờ học --- câu hỏi thường gặp ở vấn đáp.
\item Lan truyền xuôi \emph{không} tạo ma trận $BA$ kích thước $d\times k$:
\begin{quote}\small\ttfamily
return torch.add(self.base(x), linear(linear(x, self.A), self.B), alpha=self.scale)
\end{quote}
Tính $Ah$ rồi $B(Ah)$ tốn $r(d+k)$ phép nhân, so với $dk$ nếu dựng $BA$ trước.
\item Ràng buộc $1\le r\le\min(d,k)$ được kiểm ở hàm khởi tạo. Chặn trên không cần
cho tính đúng đắn của Mục~\ref{sec:b4}, nhưng $r>\min(d,k)$ không thêm biểu diễn
nào mà vẫn tốn tham số.
\end{enumerate}

Hợp nhất xuất ra một \texttt{nn.Linear} độc lập, giữ bias gốc:

\begin{quote}\small\ttfamily
merged = deepcopy(self.base)\\
merged.weight.addmm\_(self.B, self.A, alpha=self.scale)\\
return merged.eval()
\end{quote}

Hàm trả lớp mới và \emph{không} thay đổi wrapper, nên vẫn học tiếp được sau khi đã
xuất bản suy luận.

\subsubsection{Đối chiếu \dalam}

Notebook A.3 kiểm trên CPU: số tham số học đúng bằng $r(d+k)$; đầu ra của lớp đã
hợp nhất khớp đầu ra hai nhánh \emph{sau khi cập nhật khác $0$} (nếu kiểm lúc
$B$ vẫn bằng $0$ thì phép kiểm vô nghĩa, vì cả hai đều bằng lớp nền); gradient
không chảy vào weight và bias của lớp nền.

%--------------------------------------------------------------------
\subsection{A.4 --- Hàm mất mát DPO}
\label{sec:a4}

\subsubsection{Dẫn xuất}

Mục~\ref{sec:b5} chứng minh hai kết quả rồi ghép lại: nghiệm của bài toán reward
có phạt KL là phân phối Gibbs
$\pi^*\propto\pref\exp(R/\beta)$, và mô hình Bradley--Terry cho xác suất ưu tiên
là $\sigma(R_w-R_l)$. Giải nghịch quan hệ thứ nhất theo reward rồi thế vào quan hệ
thứ hai, hạng $\beta\log Z(x)$ triệt tiêu vì hai đáp án thuộc \emph{cùng} một câu
hỏi, và còn lại
\begin{equation}
  \mathcal L_{\mathrm{DPO}}(\theta)
    =-\E\bigl[\log\sigma(z_\theta)\bigr],\qquad
  z_\theta=\beta\Bigl[
    \log\tfrac{\pi_\theta(y_w\mid x)}{\pref(y_w\mid x)}
    -\log\tfrac{\pi_\theta(y_l\mid x)}{\pref(y_l\mid x)}\Bigr],
  \label{eq:dpo-a4}
\end{equation}
đúng phương trình (7) của \cite{dpo}. Không cần tính $Z(x)$ và không cần huấn
luyện reward model riêng.

\subsubsection{Cài đặt}

\begin{quote}\small\ttfamily
margin = (pw - pl) - (rw - rl)\\
return -logsigmoid(beta * margin).mean()
\end{quote}

Bốn quyết định về tính đúng đắn số học:

\begin{enumerate}[leftmargin=2em]
\item \textbf{Reference không nhận gradient:} \texttt{ref\_chosen.detach()} và
\texttt{ref\_rejected.detach()}. $\pref$ cố định theo định nghĩa bài toán; nếu
gradient chảy vào nó thì \eqref{eq:dpo-a4} không còn là mục tiêu đang tối ưu.
\item \textbf{Dùng \texttt{logsigmoid}, không dùng \texttt{log(sigmoid(\ldots))}.}
Khi $z_\theta$ âm lớn, \texttt{sigmoid} làm tròn xuống $0$ và log trả $-\infty$.
\texttt{logsigmoid} tính ổn định trên toàn miền.
\item \textbf{Không trừ hoặc nhân $\beta$ ở độ chính xác nửa:} tensor \texttt{fp16}
được nâng lên \texttt{float32} trước khi lấy hiệu. Hiệu hai log-xác suất chuỗi là
hiệu hai số lớn gần nhau --- trường hợp xấu nhất của sai số làm tròn.
\item \textbf{$\beta$ phải hữu hạn và dương:} kiểm ở đầu hàm.
Mục~\ref{sec:b5} cần $\beta>0$ ở ba chỗ (để $Z(x)$ hữu hạn, để chia $R_y/\beta$,
và để $\argmax J=\argmin\KL$), nên đây là ràng buộc toán học chứ không phải phòng
vệ đầu vào.
\end{enumerate}

Mỗi đầu vào là \emph{tổng} log-xác suất các token đáp án cho cùng prompt, không
gồm prompt và padding, và \emph{không} chia cho độ dài. Mục~\ref{sec:b6} giải thích
vì sao DPO cộng còn perplexity lại chia.

\subsubsection{Đối chiếu \dalam\ và hạng mục còn thiếu \chualam}

Notebook A.4 so mất mát và \emph{cả hai} gradient theo log-xác suất policy với
cross-entropy hai lớp trên CPU, dùng số minh họa. Bài C.4
(Mục~\ref{sec:c4}) kiểm lại trên ba dấu của margin và bốn giá trị $\beta$, khớp
autograd.

Tuy nhiên yêu cầu G5 của đề --- ít nhất một cài đặt Phần~A được dùng hoặc so sánh
\emph{trực tiếp trong quy trình Phần~D} --- \textbf{chưa} được đáp ứng: hàm này
chưa chạy trên một minibatch dữ liệu ưu tiên thật của Phần~D, vì chưa có dữ liệu
đó. Đây là cầu nối A$\leftrightarrow$D mà nhóm dự định dùng, và nó bị chặn bởi
D.3/D.4 (Mục~\ref{sec:phan-d}).

%--------------------------------------------------------------------
\subsection{Tổng hợp trạng thái Phần A}

\begin{center}
\begin{tabular}{llp{0.42\textwidth}}
\toprule
\textbf{Mã} & \textbf{Trạng thái} & \textbf{Bằng chứng}\\
\midrule
A.1a--A.1e & \dalam & \texttt{src/attention.py}; đối chiếu SDPA và MultiheadAttention\\
A.2 & \dalam & PPL $10{,}1948\pm0{,}0628$, 5 seed; chưa đối chiếu log môi trường\\
A.3 & \dalam & \texttt{src/lora.py}; đếm tham số và kiểm hợp nhất\\
A.4 & \motphan & \texttt{src/dpo.py} đã kiểm số học; chưa chạy trên minibatch thật của D\\
\bottomrule
\end{tabular}
\end{center}
